\documentclass{article}
\usepackage{graphicx}
\usepackage{amsmath}
\usepackage{amssymb}
\usepackage{amsfonts}
\usepackage{hyperref}
\usepackage{exercises}
\usepackage{indentfirst}
\usepackage{circledsteps}
\usepackage{amsthm}
\usepackage{csquotes}
\usepackage{yhmath}
\usepackage{comment}
\usepackage{placeins}
\usepackage{geometry}
\usepackage{stmaryrd}
\usepackage{hyperref}
\usepackage{tikz}
\usepackage{tabto}
\newcommand*{\oldneg}{\mathord{\sim}}
\newcommand{\Z}{\mathbb{Z}}
\newcommand{\N}{\mathbb{N}}
\let\oldemptyset\emptyset
\newcommand{\lbc}{\llbracket}
\newcommand{\rbc}{\rrbracket}

\theoremstyle{definition}
\newtheorem{definition}{Definition}
\newtheorem{theorem}{Theorem}[section]
\newtheorem{corollary}{Corollary}[theorem]
\newtheorem{lemma}[theorem]{Lemma}

\setlength{\parindent}{.5in}
\geometry{margin=0.5in}
\setlength{\parskip}{12pt}
\title{Chapter 17 Notes}
\begin{document}
\date{}
\maketitle

\section*{Section 1}

\subsection*{Notes}
Here is an explanation for the remarks given on page 224. Note that we are certainly assuming $M$ contains all ordinals for (explain why 
this is; show why this causes your given proof of the second part of remark (3) to not follow through, as it would mean that, indeed,
the sequence $R^\mathcal{G}_0,R^\mathcal{G}_1,...$ did not contain all ordinals, hence was not an ordinal hierarchy).
Definition 1.1 to even work in the first place.

\begin{lemma}
For any $\alpha$, if, $R^\mathcal{G}_\beta\in M$ for all $\beta<\alpha$, then 
both $R^\mathcal{G}_\alpha\in M$ and the $R^\mathcal{G}_\alpha$ function are in $M$. 
\end{lemma}

\begin{proof}
By transfinite induction.

Suppose it holds for all ordinals less than some ordinal $\alpha$. So suppose that 
$R^\mathcal{G}_\beta\in M$ for all $\beta<\alpha$. Then it follows that the 
$R^\mathcal{G}_\beta$ function is in $M$ for all $\beta<\alpha$. We first show that 
$R^\mathcal{G}_\alpha\in M$. We do this by showing that the 
$R^\mathcal{G}_\alpha$ function is definable over $M$, and the fact that $\alpha\in M$ (because all 
ordinals are in $M$). Then we would use condition $F_8$ of being a first-order universe 
to conclude that $R^\mathcal{G}_\alpha\in M$. Well, if $\alpha$ is a successor ordinal of some $\beta$, then
the $R^\mathcal{G}_\alpha$ function is definable over $M$ by the formula
\[ (x=\beta\wedge y=R^\mathcal{G}_\beta)\vee \]
\[ (\exists\delta<\alpha)(\exists f)(\textrm{R-delta func}(f)\wedge x\in\textrm{Dom}(f)\wedge y=f(x)) \]
and if $\alpha$ is a limit ordinal, then it is defined over $M$ by the following formula
\[ (\exists\delta<\alpha)(\exists f)(\textrm{R-delta func}(f)\wedge x\in\textrm{Dom}(f)\wedge y=f(x)) \]
Thus, in any case, we know that the $R^\mathcal{G}_\alpha$ function is definable over $M$. Thus, 
because $\alpha\in M$, we can conclude that $R^\mathcal{G}_\alpha\in M$ (as it is the domain of the function).

We first show that $R^\mathcal{G}_\alpha\in M$. Well, if $\alpha$ is a successor ordinal of some $\beta$, then it 
is defined by the formula
\[ (\forall y\in x)(y\susbseteq\mathcal{G}\times R^\mathcal{G}_\beta) \]
which works, as $\mathcal{G}\times R^\mathcal{G}_\beta\in M$ for the following reason:
both $\mathcal{G}$ and $R^\mathcal{G}_\beta$ are in $M$, meaning 
$\{\mathcal{G},R^\mathcal{G}_\beta\}\in M$ (because $M$ is a first order universe). Thus
$\cup\{\mathcal{G},R^\mathcal{G}_\beta\}\in M$, meaning
$\mathcal{P}(\cup\{\mathcal{G},R^\mathcal{G}_\beta\})\cap M\in M$, and thus
(1) $\mathcal{P}(\mathcal{P}(\cup\{\mathcal{G},R^\mathcal{G}_\beta\})\cap M)\cap M\in M$. Now, it is 
obvious that the formula in $M$ in (1) is such that $\langle x,y\rangle$ is in it for each 
$x\in\mathcal{G}$ and each $y\in R^\mathcal{G}_\beta$. Then it is easy to see that 
$\mathcal{G}\times R^\mathcal{G}_\beta$ is first-order definable over $M$, and, because this 
would mean it is a first-order definable susbset of the set in (1), 
$\mathcal{G}\times R^\mathcal{G}_\beta\in M$ (because $M$ is first-order swelled). 

The case where $\alpha$ is a limit ordinal is mostly similar. In particular, because we 
are assuming the $R^\mathcal{G}_\beta$ function is in $M$ for each $\beta<\alpha$, the following 
defines $R^\mathcal{G}_\alpha$:
\[ (\exists\beta<\alpha)(\exists f)(\textrm{R-beta func}(f)\wedge(\exists\delta<\beta)(x\in f(\delta))) \]

Thus we know that $R^\mathcal{G}_\alpha\in M$.

Now we show that the $R^\mathcal{G}_\alpha$ function is definable over $M$.
Well, if $\alpha$ is a successor ordinal of some $\beta$, then
the $R^\mathcal{G}_\alpha$ function is definable over $M$ by the formula
\[ (x=\beta\wedge y=R^\mathcal{G}_\beta)\vee \]
\[ (\exists\delta<\alpha)(\exists f)(\textrm{R-delta func}(f)\wedge x\in\textrm{Dom}(f)\wedge y=f(x)) \]
and if $\alpha$ is a limit ordinal, then it is defined over $M$ by the following formula
\[ (\exists\delta<\alpha)(\exists f)(\textrm{R-delta func}(f)\wedge x\in\textrm{Dom}(f)\wedge y=f(x)) \]
Thus, in any case, we know that the $R^\mathcal{G}_\alpha$ function is definable over $M$. Because 
$\alpha\in M$, this means that the domain of the $R^\mathcal{G}_\alpha$ function, the 
set $S$ of all $R^\mathcal{G}_\beta$ for $\beta<\alpha$, is in $M$. Hence 
$\{\alpha,S\}\in M$, $\cup\{\alpha,S\}\in M$, 
$\mathcal{P}(\cup\{\alpha,S\})\in M$,
and $\mathcal{P}(\mathcal{P}(\cup\{\alpha,S\}))\in M$. Now, 
for each $\beta<\alpha$ and $x\in S$, it of course follows that 
$\langle \beta,x\rangle\in\mathcal{P}(\mathcal{P}(\cup\{\alpha,S\}))$. That is, 
the $R^\mathcal{G}_\alpha$ function is a first-order definable subset of 
$\mathcal{P}(\mathcal{P}(\cup\{\alpha,S\}))$, which is in $M$, meaning
the function itself is in $M$.
\end{proof}  

\subsubsection*{Remark 1}
We show that each $R^\mathcal{G}_\alpha$ is a memver of $M$.

\begin{proof}
By transfinite induction. The case for 0 obviously holds (because $M$ is a first-order universe). 
Suppose the result holds for some $\alpha$. Then 
$R^\mathcal{G}_{\alpha+1}=M\cap\mathcal{P}(\mathcal{G}\times R^\mathcal{G}_\alpha)$, and thus 
$R^\mathcal{G}_{\alpha+1}\in M$ because $M$ is a first order universe. 

Now suppose the result holds for all $\alpha<\lambda$ for a limit ordinal $\lambda$.
\end{proof}

\subsubsection*{Remark 2}
We show that $R^\mathcal{G}_\alpha\subseteq R^\mathcal{G}_{\alpha+1}$. In fact, we directly prove the stronger result, 
that $\beta<\alpha$ implies $R^\mathcal{G}_\beta\subseteq R^\mathcal{G}_\alpha$.

\begin{proof}
By transfinite induction.

The least ordinal case is vacuous.

Suppose the result holds for some $\alpha$. Now, if $x\in R^\mathcal{G}_\alpha$, 
there is some $\beta<\alpha$ such that all members of $x$ are of the form 
$\{\langle y,z\rangle\}$ where $y\in\mathcal{G}$ and $z\in R^\mathcal{G}_\beta$. By the induction 
assumption, this means that $z\in R^\mathcal{G}_\alpha$. Then 
$\{\langle y,z\rangle\}$ would also be in $R^\mathcal{G}_\alpha$ for all such $y$ and $z$, and then so would 
$x$ (as we are assuming $x$ is in $M$).

The limit ordinal case is obvious.
\end{proof}

\subsubsection*{Remark 3}
We first show that $R^\mathcal{G}_{\alpha+1}\ne R_\mathcal{G}_\alpha$. For this, we prove the following lemma. Call 
a function a $R^\mathcal{G}_\alpha$ function is it has a domain of $\alpha$, and it assigns $R^\mathcal{G}_\beta$ to 
each $\beta<\alpha$. If an $R^\mathcal{G}_\alpha$ function is definable over $M$, let 
$\textrm{R-alpha func}(f)$ be a formula meaning $x$ is that function. For instance, that formula could be
\[ \textrm{Func}(f)\wedge\textrm{Dom}(f)=\alpha\wedge f(0)=0\wedge \]
\[ (\forall \beta<\alpha)[\beta^+<\alpha\supset (x\in f(\beta^+)\equiv x\subseteq\mathcal{G}\times f(\beta))] \]

\begin{lemma}
For each $\alpha$, $R^\mathcal{G}_\alpha$ is first-order definable over $M$, and 
the $R^\mathcal{G}_\alpha$ function is definable over $M$. 
\end{lemma}

\begin{proof}
By transfinite induction. Obviously the least ordinal case works.

Successor ordinal case: Suppose the result holds for some $\alpha$. Thus 
$R^\mathcal{G}_\alpha$ is definable over $M$. Also note that, by (1) of the 
remarks, $R^\mathcal{G}_\alpha\in M$. Then $R^\mathcal{G}_{\alpha+1}$ is defined over 
$M$ by the formula
\[ (\forall x\in z)(\forall y\in x)(x\in\mathcal{G}\times R^\mathcal{G}_\alpha) \] which works, 
as $\mathcal{G}$ is definable over $M$. Lastly, the $R^\mathcal{G}_{\alpha^+}$ function is definable over 
$M$ by 
\[ (x=\alpha\wedge y=R^\mathcal{G}_\alpha)\vee(\exists beta<\alpha^+)(\exists f)[ \]
\[ \textrm{R-alpha func}(f)\wedge\textrm{Ord}(x)\wedge y=f(x)] \]

Limit ordinal case: Suppose the result holds for all $\alpha<\lambda$ for some 
limit ordinal $\lambda$. Note that $\lambda\in M$ (meaning it is first-order definable over $M$, as $M$ is transitive). 
Then the following formula 
defines $R^\mathcal{G}_\lambda$ over $M$ (which works, as $M$ is assumed to be transitive):
\[ (\exists \beta<\lambda)(\exists f)([\textrm{Func}(f)\wedge\textrm{Dom}(f)=\beta\wedge f(0)=0 \]
\[ (\forall \alpha<\beta)(\alpha^+<\beta\supset (x\in f(\alpha^+)\equiv x\subseteq\mathcal{G}\times f(\alpha)))] \]
\[ \wedge (\exists \alpha<\beta)(x\in f(\alpha))) \]
Lastly, the $R^\mathcal{G}_{\lambda}$ function is definable over $M$ by
\[ (\exists \alpha<\lambda)(\textrm{R-alpha func}(f)\wedge\textrm{Ord}(x)\wedge x=f(y) \]
\end{proof}

Note that, as a consequence of the above lemma (and part (2) of the remarks), 
we know that, for all $\alpha$, $R^\mathcal{G}_\alpha\in R^\mathcal{G}_{\alpha+1}$. Now we show that 
$R^\mathcal{G}_{\alpha}\ne R^\mathcal{G}_{\alpha+1}$.

\begin{proof}
This is in fact immediate from the above lemma. For, we know from that that 
$R^\mathcal{G}_\alpha\in R^\mathcal{G}_{\alpha+1}$. But this means, if 
$R^\mathcal{G}_\alpha=R^\mathcal{G}_{\alpha+1}$, then 
$R^\mathcal{G}_\alpha\in R^\mathcal{G}_\alpha$. By (1) of the remarks, we know that 
$R^\mathcal{G}_\alpha$ is a subclass of $M$, hence well founded (because we are assuming $M$ to 
be well founded). But $R^\mathcal{G}_\alpha\in R^\mathcal{G}_\alpha$ contradicts this.

By transfinite induction. In particular, we show that, for every $\alpha$, $R^\mathcal{G}_\beta\ne R^\mathcal{G}_\alpha$ for every 
$\beta<\alpha$. We do this by assuming the result holds for all ordinals l

Suppose the result holds for all ordinals less than some ordinal $\delta$.

Successor ordinal case: Suppose the $\delta$ is the successor of some ordinal $\alpha$.
There are two cases: either $\alpha$ itself is a successor ordinal, 
or it is a limit ordinal (of course, $\alpha$ could be $\emptyset$, but that case is trivial). We must show the result holds for $\alpha^+$. 
Note that it suffices to show that $R^\mathcal{G}_{\alpha+1}\ne R^\mathcal{G}_\alpha$, for, if we show that, then if 
$R^\mathcal{G}_{\alpha+1}=R^\mathcal{G}_\beta$ for some $\beta<\alpha$, then it would follow that 
$R^\mathcal{G}_\beta=R^\mathcal{G}_\alpha$ by part (2) of the remarks, contrary to the induction assumption.

On the one hand, suppose that $\alpha$ itself is a successor ordinal (say of some $\beta$). Thus, by the assumption, we 
know that there must be some $y\in R^\mathcal{G}_\alpha$ such that $y\notin R^\mathcal{G}_\beta$. Consider any $x\in\mathcal{G}$. 
Then $\{\langle x,y\rangle\}\in R^\mathcal{G}_{\alpha+1}$ (because $M$ is a first-order universe). If 
$R^\mathcal{G}_{\alpha+1}=R^\mathcal{G}_\alpha$, then it would follow that $z=\{\langle x,y\rangle\}$ for some 
$z\in R^\mathcal{G}_\alpha$. What this would mean is that $\langle x,y\rangle\in\mathcal{G}\times R^\mathcal{G}_\beta$. 
But that would mean $y\in R^\mathcal{G}_\beta$, a contradiction. Thus it must be that $R^\mathcal{G}_{\alpha+1}\ne R^\mathcal{G}_\alpha$.

On the other hand, suppose that $\alpha$ is a limit ordinal. Then pick some successor ordinal $\beta^+<\alpha$. By the 
transfinite induction assumption, we know that there is some $x\in R^\mathcal{G}_{\beta^+}$ such that 
$x\notin R^\mathcal{G}_\beta$. By (2) of the remarks, we know that this means $x\in R^\mathcal{G}_
\end{proof}

Now we show that $\mathcal{D}^\mathcal{G}$ is a proper class in $M$ (that is, 
$\mathcal{D}^\mathcal{G}\notin M$). Suppose it were in $M$. Then, because 
$R^\mathcal{G}_0,R^\mathcal{G}_1,...,R^\mathcal{G}_\alpha,...$ is an ordinal hierarchy, it 
follows by Theorem 2.1 of Chapter 7 that there is some ordinal $\alpha$ such that 
$\mathcal{D}^\mathcal{G}\subseteq R^\mathcal{G}_\alpha$. By (2), 
$\mathcal{D}^\mathcal{G}\subseteq R^\mathcal{G}_{\alpha+1}$.
Also $R^\mathcal{G}_{\alpha+1}\subseteq\mathcal{D}^\mathcal{G}$, meaning 
$R^\mathcal{G}_{\alpha+1}=\mathcal{D}^\mathcal{G}$. But this contradicts the fact that
$R^\mathcal{G}_{\alpha+1}\ne R^\mathcal{G}_\alpha$. Thus it cannot be that 
$\mathcal{D}^\mathcal{G}\in M$.

\subsubsection*{Remark 4}
Suppose $f\subseteq\mathcal{G}\times\mathcal{D}^\mathcal{G}$ and $f$ is an $M-set$. Now, for 
each $x\in f$, $x\in R^\mathcal{G}_\alpha$ for some $\alpha$. Let $o(f)$ be the class of all 
$\alpha$ such that $x\in f$ for at least one $x\in R^\mathcal{G}_\alpha$. We show that there must be some 
$\beta$ higher than all the $\alpha\in o(f)$. For suppose there weren't. Then it would follow that 
$o(f)$ could be put into a 1-1 correspondence with $On$. But $o(f)$ is a set (because $f$ is a set), which would 
mean that $On$ is a set (by the axiom of substitution), which is impossible. Thus it must be that there is some 
$\beta$ higher than all the $\alpha\in o(f)$. Consider the least such $\beta$. Because $f\in M$, it is easy to see that 
$f\in R^\mathcal{G}_{\beta+1}$, and thus $f\in\mathcal{D}^\mathcal{G}$.

\subsection*{Remark 5}
I could not see how to show the isomorphism, but the rest of the remark seems to be just Exercise 1.2, which I will give here.

Here is how we could get this recursion to work in the first place. We define a function $m_\alpha$ for each 
$\alpha$ whose domain is $R^\mathcal{G}_\alpha$ and which satisfies the conditions of the function. We use 
transfinite recursion, where $m_{\alpha+1}=g(m_\alpha)$, where 
$g(m_\alpha)$ is the set containing all members of $m_\alpha$, as well as all 
sets of the form $\langle x,\{m_\alpha(g):\langle p,g\rangle\in x\}$. 

Suppose $\mathcal{G}=\{p\}$, and a mapping $m$ is defined by (I was unsure how we could get this recursion to work in the first place)
\[ m(f)=\{m(g)|\langle p,g\rangle\in f\} \]

First, we show that the mapping is 1-1. Note for this that an obvious transfinite induction argument shows that
if $\beta<\alpha$, $m(g)$ stays constant between $m$ applied to $R^\mathcal{G}_\beta$ and $R^\mathcal{G}_\alpha$, for 
any $m\in R^\mathcal{G}_\beta$.

\begin{proof}
By transfinite induction.

The case for $R^\mathcal{G}_0$ obviously holds.

Suppose it holds for all ordinals less than some $\alpha$. Now consider any 
$f$ and $h$ in $R^\mathcal{G}_\alpha$. Suppose $m(f)=m(h)$. Then 
\[ m(f)=\{m(g)|\langle p,g\rangle\in f\}=m(h)=\{m(l)|langle p,l\rangle\in f\} \]
Then it easily follows that $f=h$, as the values of $m(g)$ for all 
$g\in R^\mathcal{G}_\beta$ for some $\beta<\alpha$ stay constant. 
\end{proof}

Next we show that, for all $f\in R^\mathcal{G}_\alpha$, $m(f)\in R_\alpha$.

\begin{proof}
By transfinite induction.

The case for $R^\mathcal{G}_0$ vacuously holds.

Suppose it holds for some $\alpha$. Consider any $x\in R^\mathcal{G}_{\alpha+1}$; then
$m(f)=\{m(g)|\langle p,g\rangle\in f\}$, and, by assumption, each 
$m(g)$ would be in $R_\alpha$. Thus $m(f)$ is a subset of $R_\alpha$, and, consequently, 
it is in $R_{\alpha+1}$ (because $R_{\alpha+1}=\mathcal{P}(R_\alpha)}$.
\end{proof}

\end{document}
